%!TEX TS-program = pdflatex
%!TEX encoding = UTF-8 Unicode
% =============================================================
%   Proyecto Final de Ingeniería Informática — USAL
%   Archivo principal. Compilar con pdfLaTeX (+ biber + makeglossaries).
%   En Overleaf: Menú > Compiler > pdfLaTeX. La primera compilación del
%   plan gratuito puede cortarse a los 20 s: apretar Recompile otra vez
%   y termina (queda caché). Después cada compilación tarda unos segundos.
% =============================================================

\documentclass[%
    language=spanish,   % spanish (por defecto) o english
    guias=true,         % true: muestra las cajas de guía de la cátedra
                        % false: las oculta (usar para la ENTREGA FINAL)
    pdfa=false,         % true: genera PDF/A-2b (más lento; requiere LuaLaTeX)
]{usal-proyecto-final}


% -------------------------------------------------------------
% DATOS DEL TRABAJO  (completar)
% -------------------------------------------------------------
\title{Tinku: una plataforma digital de marketplace educativo para\\
       conectar estudiantes y tutores académicos a escala nacional}
\thesisHeader{Facultad de Ingeniería\\Proyecto Final de Ingeniería Informática}
\academicYear{2026}
% \graduationDate{17/11/2026} % (opcional) fecha de defensa

% autor/a (para trabajos grupales, separar con \\ o \and)
\author{Santiago Mercado Carbone}
\authorEmail{s.mercadocarbone@usal.edu.ar}

% autoridades que figuran en la portada (verificar cargos vigentes)
\decano{Walter Oscar Rodríguez Esquivel}{Decano, Facultad de Ingeniería}
\director{Germán Buchniv}{Director de las carreras Lic. en Sistemas de Información e Ingeniería en Informática}
\profesor{Esteban Tissera}{Profesor, Facultad de Ingeniería}
% \tutor{Nombre del tutor o tutora}{Cargo, Institución} % (opcional)

% resumen / abstract
\palabrasclave{marketplace educativo, tutorías particulares, inteligencia artificial, matching semántico, educación personalizada}
\keywords{educational marketplace, private tutoring, artificial intelligence, semantic matching, personalized education}

% segunda página: declaración de originalidad
\secondpage{%!TEX root = ../main.tex
% Segunda página (reverso de la portada): declaración de originalidad.

\mbox{}
\vfill

{\small\onehalfspacing
\begin{center}

\noindent\rule{\textwidth}{1pt} % ----------

    Este trabajo es original y fue elaborado exclusivamente con este fin.
    Todos los autores cuyos estudios contribuyeron a su elaboración se
    encuentran debidamente citados. Se permite su reproducción parcial con
    indicación del autor, del grado, del año académico, de la institución
    (Universidad del Salvador) y de la fecha de defensa pública.

\noindent\rule{\textwidth}{.2pt} % ----------

\end{center}
} % ----------
}

% dedicatoria (opcional; descomentar para agregar una página de dedicatoria)
% \dedication{\textit{A ...}}

% portada: logo y marca de agua (opcionales)
% \coverlogo{template/usal}            % logo superior (por defecto)
% \coverwatermark{template/marca-agua} % marca de agua de fondo (vacío = sin marca)


% -------------------------------------------------------------
% RECURSOS
% -------------------------------------------------------------
\addbibresource{refs.bib}                 % bibliografía (BibTeX)
\loadglsentries{preliminares/glosario}    % términos del glosario
\loadglsentries{preliminares/siglas}      % siglas y acrónimos
\loadglsentries{preliminares/simbolos}    % símbolos
%!TEX root = ./main.tex
% Macros, teoremas y atajos propios del trabajo.

% entornos tipo teorema, numerados por capítulo
\newtheorem{definicion}{Definición}[chapter]
\newtheorem{proposicion}{Proposición}[chapter]
\newtheorem{teorema}{Teorema}[chapter]

% nombre del sistema o producto, para no repetirlo a mano
\newcommand{\sistema}{\textsc{Tinku}}

% atajos matemáticos de ejemplo
\newcommand{\vect}[1]{\mathbf{#1}}
\DeclareMathOperator{\argmax}{arg\,max}
\DeclareMathOperator{\argmin}{arg\,min}

% moneda: \usd{77200} -> USD 77.200
\newcommand{\usd}[1]{USD~#1}

% comisión/referencia del ticket promedio de sesión (Cap. 5)
\newcommand{\comision}{27\,\%}
                    % macros y teoremas propios

% índice general: 2 = hasta subsección (1.1.1 no se lista)
\setcounter{tocdepth}{2}

% evita títulos de sección huérfanos al pie de página
\RequirePackage{needspace}
\RequirePackage{etoolbox}
\pretocmd{\section}{\Needspace{12\baselineskip}}{}{}
\pretocmd{\subsection}{\Needspace{9\baselineskip}}{}{}
\pretocmd{\subsubsection}{\Needspace{5\baselineskip}}{}{}

%% Para imprimir solo la tapa con letra grande y sin marca de agua:
%\begin{document}
%    \makecover
%\end{document}


% =============================================================
\begin{document}

    % las páginas no se estiran para llenar (evita saltos de párrafo largos)
    \raggedbottom

    % portada, agradecimientos, resumen, declaración de IA, índices
    \frontmatter

    % cuerpo del trabajo: un archivo por capítulo, en el orden de entrega
    \mainmatter
    %!TEX root = ../main.tex
% =============================================================
%  CAPÍTULO 1 — INTRODUCCIÓN            (Entrega 1 · 5 % · 9/6)
% =============================================================

\chapter{Introducción}\label{cap:introduccion}

\begin{guia}[Guía de la cátedra: qué se evalúa en este capítulo]
\begin{enumerate}
    \item \textbf{Descripción del problema}: concreto, anclado en evidencia
          (datos duros con fuente, casos reales, competidores nombrados).
    \item \textbf{Objetivo general}: uno solo, medible, en un párrafo.
          Los \textbf{objetivos específicos} se desprenden de él y
          \emph{no} son los pasos del plan de trabajo.
    \item \textbf{Hipótesis}: explícita, formulada como pregunta o
          afirmación de investigación que pueda responderse en el
          capítulo~\ref{cap:conclusiones}.
    \item \textbf{Relevancia}: por qué importa, a quién beneficia, qué les
          falta a las soluciones existentes.
    \item \textbf{Alcance y recursos}: qué entra y qué no; stack y fuentes
          de datos previstos. El ``alcance que nunca cierra'' es el riesgo
          número uno de la materia.
\end{enumerate}
Extensión orientativa: 4 a 6 páginas. Redacción impersonal, sin ``yo
creo'', oraciones de 15 a 22 palabras, sin palabras vacías
(``claramente'', ``obviamente''). Toda cifra con su fuente citada, por
ejemplo \citep{breiman2001}.
\end{guia}

\completar{Desarrollo del capítulo. Las secciones y subsecciones se
organizan como resulte mejor para el proyecto.}

\section{\completar{Nombre de la sección}}

\begin{ejemplo}[Ejemplo: forma de un objetivo general]
Desarrollar un sistema de \completar{qué} para \completar{quién} en
\completar{dónde}, que permita \completar{resultado medible} mediante
\completar{enfoque o tecnología central}.
\end{ejemplo}

\completar{Texto.}

\subsection{\completar{Nombre de la subsección}}

\completar{Texto.}

    \glsresetall % las siglas se vuelven a expandir completas a partir de acá
    %!TEX root = ../main.tex
% =============================================================
%  CAPÍTULO 2 — METODOLOGÍA             (Entrega 2 · 5 % · 5/5)
% =============================================================

\chapter{Metodología y procedimientos}\label{cap:metodologia}

\begin{guia}[Guía de la cátedra: qué se evalúa en este capítulo]
La cátedra exige \textbf{dos metodologías}. El error más común es
entregar solo una.
\begin{enumerate}
    \item \textbf{Metodología de la investigación}: área temática,
          planteamiento y delimitación del problema, marco teórico
          previsto, diseño concreto (fases), indicadores medibles con
          fuente y umbral, técnicas de recolección, instrumentos, datos,
          procesamiento, análisis crítico y síntesis. Tipo de
          investigación: en general, aplicada.
    \item \textbf{Metodología de dirección del proyecto}: marco de trabajo
          justificado, alcance y entregables, cronograma con fechas reales
          (Gantt con hitos que coinciden con las entregas), costos, riesgos
          (referencia al capítulo~\ref{cap:riesgos}), calidad, recursos y
          satisfacción de las partes interesadas.
    \item \textbf{Declaración de uso de IA}: qué tareas fueron asistidas y
          cómo se verificaron. Los prompts se guardan en una carpeta
          compartida con la cátedra.
\end{enumerate}
Extensión orientativa: 6 a 10 páginas. Referencias habituales:
\citet{hernandezsampieri2014} para investigación y \citet{pmi2021}
para gestión.
\end{guia}

\section{Introducción}

El presente capítulo describe el marco metodológico adoptado para el
desarrollo de \sistema{}. En línea con los requisitos del proyecto
final, la metodología se articula en dos dimensiones complementarias:
la metodología de investigación, que define cómo se aborda el problema,
se recolecta la información y se validan los resultados; y la
metodología de gestión del proyecto, que establece cómo se organiza,
planifica y ejecuta el trabajo de desarrollo a lo largo del tiempo
disponible. Ambas dimensiones son necesarias y se retroalimentan: el
proceso investigativo informa las decisiones de diseño técnico, mientras
que el marco de gestión garantiza que esas decisiones se implementen de
manera ordenada, dentro del alcance y los plazos establecidos.

\section{Tipo de investigación}

El presente proyecto se enmarca dentro de la \textbf{investigación
aplicada}. A diferencia de la investigación básica --orientada a la
producción de conocimiento teórico sin aplicación inmediata-- la
investigación aplicada parte de un problema real y concreto, y genera
conocimiento con utilidad directa para resolverlo. En el caso de
\sistema{}, el problema es la fragmentación e inaccesibilidad del
mercado de tutorías académicas en Argentina, y el conocimiento generado
se materializa en un artefacto funcional: una plataforma digital que
conecta tutores y alumnos a escala nacional. El aporte original del
proyecto reside en la aplicación de técnicas de \gls{matching} semántico
basadas en \gls{ia} a un dominio educativo local, y en la integración de
funcionalidades de tutoría virtual en un único ecosistema diseñado para
las condiciones geográficas y de conectividad del territorio argentino.
El enfoque es de carácter mixto: predominantemente cualitativo en la
fase de relevamiento del problema y cuantitativo en la fase de
validación.

\section{Metodología de la investigación}

La investigación se estructura siguiendo las etapas del proceso
científico aplicado al desarrollo de sistemas de información.

\subsection{Área temática}

El área temática del proyecto se sitúa en la intersección entre el
desarrollo de plataformas digitales de tipo \gls{marketplace}, la \gls{ia}
aplicada al \gls{matching} de usuarios, y la tecnología educativa
(\gls{edtech}).

\subsection{Planteamiento del problema}

El problema de investigación se formula como la ausencia de una
plataforma nacional que centralice, estandarice y amplíe
geográficamente el acceso al mercado de tutorías particulares en
Argentina.

\subsection{Delimitación de la investigación}

La investigación se delimita a los objetivos específicos OE1 a OE8
definidos en el capítulo~\ref{cap:introduccion}, con foco en el diseño e
implementación del \gls{mvp} de \sistema{} para el mercado argentino.
Quedan fuera del alcance investigativo los mercados regionales, la
integración con organismos estatales y los módulos de accesibilidad y
suscripción premium, los cuales se documentan como trabajo futuro.

\subsection{Marco teórico}

Se realizará una búsqueda sistemática de bibliografía académica y
técnica sobre los temas centrales del proyecto, abarcando plataformas de
\gls{marketplace} educativo, algoritmos de recomendación basados en
\gls{embeddings} semánticos, sistemas de reputación en plataformas
digitales, tecnología \gls{webrtc} aplicada a educación, y el contexto
educativo argentino. Esta revisión se documenta en el
capítulo~\ref{cap:marco-teorico}. Las fuentes a consultar incluyen
artículos científicos, tesis de grado y posgrado, documentación técnica
oficial, informes de mercado de organizaciones como Grand View Research
y Research and Markets, y reportes del Ministerio de Educación de la
Nación Argentina.

\subsection{Diseño concreto de investigación}

A partir de la información relevada, se definirá el diseño técnico del
sistema, abarcando la arquitectura, el modelo de datos, los flujos de
usuario y los contratos de \gls{api}. Esta etapa se documenta en el
capítulo~\ref{cap:marco-tecnologico}.

\subsection{Indicadores}

Los indicadores principales para verificar la hipótesis central se
resumen en la tabla~\ref{tab:indicadores}.

\begin{table}[H]
    \centering
    \caption{Indicadores del proyecto}
    \label{tab:indicadores}
    \small
    \begin{tabularx}{\textwidth}{l L L l}
        \toprule
        \tb{ID} & \tb{Variable} & \tb{Indicador} & \tb{Umbral} \\
        \midrule
        I-01 & Aceptación del \gls{matching} & Tasa de aceptación de las sugerencias del motor por parte del alumno & $> 60\,\%$ de recomendaciones aceptadas en el piloto \\
        I-02 & Calidad percibida de la sesión & Puntuación promedio otorgada por los alumnos en el cuestionario post-sesión (escala 1 a 5) & $\geq 4{,}0$ \\
        I-03 & Cobertura geográfica del piloto & Sesiones completadas entre usuarios de provincias distintas & Al menos 1 sesión exitosa entre dos provincias, incluyendo una del interior \\
        I-04 & Satisfacción y recomendación & \gls{nps} del piloto & Porcentaje positivo (promotores $>$ detractores) \\
        \bottomrule
    \end{tabularx}
\end{table}

\subsection{Técnicas de recolección de datos}

Se implementarán técnicas cualitativas y cuantitativas para la
recolección de datos primarios. La técnica cualitativa consistirá en
entrevistas para validar el problema planteado, identificar puntos de
dolor en el proceso de búsqueda y contratación, y relevar expectativas.
La técnica cuantitativa se basará en encuestas de validación de demanda
para cuantificar la frecuencia de uso de tutorías, los canales actuales
de búsqueda y la disposición a utilizar y pagar por una plataforma
centralizada.

\subsection{Instrumentos de recolección de datos}

Se utilizarán entrevistas semiestructuradas aplicadas a una muestra de
al menos dos tutores particulares activos y cuatro alumnos o familias
que hayan contratado tutorías, complementadas con un formulario breve en
Google Forms orientado a estudiantes universitarios y familias con hijos
en edad escolar, buscando un mínimo de 30 respuestas válidas. El guion
de las entrevistas se transcribe en el apéndice~\ref{apx:entrevista} y
la encuesta completa en el apéndice~\ref{apx:encuesta}.

\subsection{Datos}

Los datos a obtener consistirán en registros de audio provenientes de
las entrevistas semiestructuradas y datos tabulares generados a partir
de las respuestas numéricas y de selección múltiple del formulario
digital.

\subsection{Procesamiento de los datos}

Las entrevistas serán grabadas con el consentimiento del entrevistado
para luego ser transcritas en su totalidad y analizadas mediante
codificación temática, mientras que los datos cuantitativos obtenidos de
los formularios serán agrupados y tabulados para permitir un conteo y
evaluación estadística de las métricas.

\subsection{Análisis de los datos}

Los datos recolectados durante el piloto funcional y la validación de
demanda se analizarán en relación con los indicadores definidos,
presentando los resultados estructurados en el
capítulo~\ref{cap:resultados}.

\subsection{Síntesis y conclusiones}

La síntesis de los datos contrastados con los indicadores de validación
evaluará el cumplimiento de la hipótesis central, presentando las
conclusiones finales y las implicancias del proyecto en el
capítulo~\ref{cap:conclusiones}.

\section{Metodología de gestión del proyecto: Scrum adaptado}

Para la planificación y ejecución del desarrollo de \sistema{} se adopta
\textbf{Scrum adaptado a un contexto unipersonal}, siendo el marco ágil
más apropiado dado que se cuenta con un único desarrollador, un conjunto
de módulos técnicos bien delimitados y una fecha de entrega final fija.
La metodología en cascada fue descartada porque asume que todos los
requisitos pueden definirse con precisión al inicio y que las etapas son
completamente secuenciales, lo cual choca con la necesidad de iteración
en proyectos con \gls{ia} e integraciones complejas como LiveKit o
\gls{api} externas. Kanban fue descartado como marco principal porque,
aunque útil para tareas diarias, no provee la planificación por
objetivos ni los hitos verificables exigidos por el cronograma
académico.

Scrum organiza el trabajo en \textbf{sprints de dos semanas} con
objetivos concretos, revisión del incremento al final de cada ciclo y un
backlog priorizado. En este contexto unipersonal, el desarrollador
asume los roles de Product Owner, Scrum Master y Development Team,
contando con el director del proyecto como stakeholder externo para la
revisión académica. Los artefactos incluyen el Product Backlog, el
Sprint Backlog y el Incremento, mientras que las ceremonias adaptadas
comprenden la Sprint Planning de una hora al inicio de cada ciclo, un
Daily Stand-up diario personal documentado en la bitácora de desarrollo,
la Sprint Review y la Sprint Retrospective.

\section{Gestión del proyecto: dimensiones clave}\label{sec:dimensiones}

En línea con los criterios de evaluación del proyecto y el perfil
directivo en Project Management Professional, se detallan las
dimensiones principales de gestión:

\begin{itemize}
    \item \textbf{Alcance:} definido por los nueve módulos del \gls{mvp}
          descritos en el capítulo~\ref{cap:introduccion}. Requiere
          aprobación explícita cualquier funcionalidad adicional para el
          control del alcance, apoyándose en el uso del Product Backlog,
          el Sprint Backlog y los incrementos funcionales.
    \item \textbf{Tiempos:} el proyecto se estructura en sprints de dos
          semanas a lo largo de seis meses entre junio y noviembre de
          2026, con una dedicación estimada de 15 a 20 horas semanales e
          hitos académicos inamovibles, integrando las ceremonias del
          marco ágil adoptado.
    \item \textbf{Costos:} los costos operativos del \gls{mvp} se estiman
          entre cero y cien dólares mensuales durante el desarrollo,
          cubiertos mediante planes gratuitos y de bajo costo en la nube,
          detallándose plenamente en el capítulo~\ref{cap:economica}.
    \item \textbf{Riesgos:} el análisis detallado se presenta en el
          capítulo~\ref{cap:riesgos}, contemplando factores como demoras
          en la integración de \gls{api} externas, costos de modelos de
          lenguaje superiores a los estimados y la dificultad para
          reclutar usuarios para el piloto.
    \item \textbf{Calidad:} cada módulo desarrollado debe superar pruebas
          unitarias con una cobertura mínima del 70\,\%, además de
          sortear una batería de pruebas de integración previas al
          piloto, midiendo la calidad percibida mediante los indicadores
          específicos de la investigación.
    \item \textbf{Recursos:} un desarrollador full-stack con
          conocimientos en React, Java con Spring Boot, Python y bases de
          datos relacionales, junto al uso de herramientas tecnológicas de
          soporte y la estructura de roles unipersonales descrita.
    \item \textbf{Satisfacción de los interesados:} se medirá a través de
          un sistema de calificaciones post-clase bidireccional para
          mantener altos estándares de calidad, complementándose con el
          análisis del \gls{nps} del grupo piloto para comprobar el grado
          de utilidad percibida y la aplicabilidad de los resultados en
          escenarios reales.
\end{itemize}

\section{Cronograma de entregables}

El cronograma fue construido de atrás hacia adelante tomando como fecha
límite el 28 de noviembre de 2026, fecha estimada de la presentación
final, y distribuyendo los entregables académicos y los sprints de
desarrollo de manera que cada capítulo se apoye en trabajo ya realizado
(tabla~\ref{tab:hitos}).

\begin{table}[H]
    \centering
    \caption{Cronograma de entregables y sprints asociados}
    \label{tab:hitos}
    \small
    \begin{tabularx}{\textwidth}{L L L}
        \toprule
        \tb{Hasta} & \tb{Entregable académico} & \tb{Sprints de desarrollo asociados} \\
        \midrule
        9 jun 2026 & Cap.~2 Metodología & Sprint 0: setup del entorno, estructura del repositorio, wireframes iniciales + spike/prototipo técnico (kill-switch) \\
        28 jul 2026 & Cap.~3 Marco teórico & Sprints 1--2: autenticación y perfiles de usuario (backend + frontend) \\
        18 ago 2026 & Cap.~4 Marco tecnológico & Sprints 3--4: motor de \gls{matching} IA + sistema de búsqueda y filtros \\
        8 sep 2026 & Cap.~5 Justificación económica & Sprints 5--6: reservas + integración MercadoPago \\
        22 sep 2026 & Cap.~6 Análisis de riesgos & Sprint 7: aula virtual (LiveKit + Redis) \\
        13 oct 2026 & Cap.~7 Resultados & Sprints 8--9: módulo IA educativa + sistema de reputación \\
        27 oct 2026 & Cap.~8 Conclusiones & Sprint 10: integración end-to-end, pruebas de sistema y correcciones \\
        28 nov 2026 & Presentación final + producto funcional & Sprint 11: piloto con usuarios reales, análisis de resultados, documentación final \\
        \bottomrule
    \end{tabularx}
\end{table}

Cabe aclarar que Redis en el Sprint 7 cubre el estado propio de la
aplicación --chat en tiempo real, límite de tasa de peticiones-- y no el
estado interno de LiveKit, que al usarse en su modalidad Cloud queda
gestionado del lado del proveedor sin intervención del backend. La
figura~\ref{fig:gantt} muestra la línea de tiempo completa.

\begin{figure}[H]
    \centering
    \begin{ganttchart}[
        hgrid, vgrid={*{6}{draw=none}, dotted},
        x unit=0.42mm,
        y unit chart=6mm, y unit title=7mm,
        time slot format=isodate,
        bar label font=\small, milestone label font=\small,
        title label font=\small,
        bar/.append style={fill=mainColor!60, draw=none},
        milestone/.append style={fill=mainColor, draw=none},
    ]{2026-06-01}{2026-11-28}
        \gantttitlecalendar{month=shortname} \\
        \ganttbar{Sprint 0 --- Setup}{2026-06-01}{2026-06-09} \\
        \ganttbar{Sprints 1--2 --- Autenticación y perfiles}{2026-06-09}{2026-07-28} \\
        \ganttbar{Sprints 3--4 --- Matching IA}{2026-07-28}{2026-08-18} \\
        \ganttbar{Sprints 5--6 --- Reservas y pagos}{2026-08-18}{2026-09-08} \\
        \ganttbar{Sprint 7 --- Aula virtual}{2026-09-08}{2026-09-22} \\
        \ganttbar{Sprints 8--9 --- IA educativa y reputación}{2026-09-22}{2026-10-13} \\
        \ganttbar{Sprint 10 --- Integración}{2026-10-13}{2026-10-27} \\
        \ganttbar{Sprint 11 --- Piloto}{2026-10-27}{2026-11-27} \\
        \ganttmilestone{Entrega final}{2026-11-28}
    \end{ganttchart}
    \caption{Cronograma de sprints del proyecto (ciclo lectivo 2026)}
    \label{fig:gantt}
\end{figure}

\section{Herramientas de gestión y seguimiento}

Para garantizar la trazabilidad del trabajo y la calidad del proceso, se
implementarán diversas herramientas tecnológicas de soporte. Se
utilizará \textbf{GitHub Projects} para la gestión del backlog y los
sprints mediante tableros Kanban integrados al repositorio, donde cada
issue representa una tarea del sprint backlog. El control de versiones
se efectuará en GitHub mediante ramas de características y solicitudes
de revisión antes de la integración a la rama principal. Asimismo, se
mantendrá una bitácora de desarrollo en formato de texto simple
actualizada al cierre de cada jornada para registrar avances y decisiones
técnicas como parte del seguimiento diario.

Para el diseño de interfaz se empleará Figma en la elaboración de
prototipos de alta fidelidad, mientras que Postman servirá para la
documentación y prueba de los endpoints de la \gls{api} REST, y Jest se
configurará como marco de pruebas unitarias del frontend, mientras que
JUnit 5 junto con Mockito cumplen ese rol en el backend Spring Boot,
ambos integrados en el pipeline de integración continua mediante GitHub
Actions. Para el despliegue, el backend se empaqueta como imagen Docker
y se publica mediante Coolify, una plataforma de despliegue autoalojada
que gestiona contenedores, certificados y variables de entorno de forma
equivalente a un PaaS gestionado, alojada primero en infraestructura
propia del desarrollador y migrada a un \gls{vps} antes del piloto.
    %!TEX root = ../main.tex
% =============================================================
%  CAPÍTULO 3 — MARCO TEÓRICO Y ESTADO DEL ARTE
%                                       (Entrega 3 · 10 % · 28/7)
% =============================================================

\chapter{Síntesis de la literatura consultada}\label{cap:marco-teorico}

\begin{guia}[Guía de la cátedra: qué se evalúa en este capítulo]
\begin{enumerate}
    \item \textbf{Estado del arte}: investigaciones nacionales,
          internacionales y productos similares, por separado. De cada
          trabajo: qué hizo y qué resultados obtuvo, con métricas.
    \item \textbf{Brechas identificadas}: qué falta en lo existente y
          justifica \emph{este} proyecto. Es lo que separa un estado del
          arte de un listado. Incluir también los resultados negativos de
          la literatura.
    \item \textbf{Bases teóricas}: desarrollan las variables del título y su
          relación. Cada concepto termina explicando cómo se usa en el
          proyecto (teoría instrumental, no de relleno). Citas canónicas
          antes que tutoriales.
    \item \textbf{Marco conceptual}: conceptos que se nombran pero no se
          desarrollan. Los términos del glosario se marcan con
          \verb|\gls{}|, por ejemplo \gls{g:mvp} y \gls{prototipo}.
    \item \textbf{Calidad bibliográfica}: papers, tesis y libros de
          editoriales serias, con datos completos. Wikipedia y blogs restan.
    \item \textbf{Entrevistas a especialistas} (si las hay): fecha, nombre,
          cargo y aporte; transcripción en el apéndice~\ref{apx:entrevista}.
\end{enumerate}
Errores que la cátedra marca: falta de coherencia interna, bibliografía
superficial, información irrelevante o redundante. Es uno de los
capítulos más extensos (orientativo: 20 a 30 páginas).
\end{guia}

\section{Estado del arte}\label{sec:estado-del-arte}

\subsection{Investigaciones nacionales}

\citet{torino2023}, en su tesis de maestría de la Universidad Torcuato
Di Tella, desarrolló un modelo predictivo de deserción escolar para
Argentina utilizando datos socioeconómicos y habitacionales de la
Encuesta Permanente de Hogares, alcanzando un desempeño del 87,4\,\%
bajo la métrica AUC-ROC. Un hallazgo relevante para este proyecto es que
la autora señala la ausencia, a nivel nacional, de un sistema de
seguimiento educativo que permita identificar tempranamente a los
estudiantes en riesgo para brindarles un acompañamiento personalizado
--la misma carencia estructural de personalización que \sistema{} busca
atender desde el lado de la oferta de tutorías.

Un estudio sobre ruralidad y educación en Argentina \citep{olea}
cuantifica la brecha urbano-rural en la finalización del nivel
secundario: mientras que en Tierra del Fuego la diferencia es inferior a
10 puntos porcentuales, en provincias del norte del país --Chaco,
Corrientes, Misiones, Santiago del Estero, Jujuy y Tucumán-- la brecha
supera los 30 puntos porcentuales entre estudiantes urbanos y rurales de
15 a 17 años. El mismo trabajo documenta que algunas jurisdicciones ya
recurrieron históricamente a esquemas de ``escuelas mediadas por
tecnología con docentes en las ciudades y estudiantes en sus localidades
acompañados por tutores'' como estrategia para ampliar el acceso --un
antecedente conceptual directo del modelo remoto que propone \sistema{}.

Ambos antecedentes empíricos evidencian, por un lado, la necesidad
crítica de implementar sistemas de seguimiento personalizado y, por el
otro, validan estructuralmente la eficacia de los modelos mediados por
tecnología para sortear las barreras de dispersión geográfica en el
territorio nacional.

\subsection{Investigaciones internacionales}

Respecto del impacto de la tutoría personalizada, una revisión
sistemática con meta-análisis de tres niveles \citep{sciencedirect2022}
sintetizó evidencia experimental sobre tutoría privada y encontró
tamaños de efecto de entre 0,42 y 0,67 sobre el rendimiento académico,
con la región y la materia como las variables moderadoras más
influyentes entre los trece factores analizados. En la misma dirección,
un trabajo reciente sobre tutoría bajo demanda \citep{ondemand2026}
sostiene que la tutoría es una de las intervenciones educativas más
efectivas para mejorar los resultados de aprendizaje, superando en
muchos casos a otras intervenciones como la reducción del tamaño de
clase. Ambos antecedentes respaldan empíricamente la Variable~2 del
marco teórico de este proyecto (sección~\ref{sec:variable2}).

En el plano tecnológico, un trabajo de la Universidad de Ottawa
\citep{isotropy2026} propuso un sistema de recomendación semántica de
cursos universitarios basado en aprendizaje contrastivo sobre
representaciones BERT, aplicado a un catálogo de más de 500 cursos de
ingeniería. El antecedente es relevante porque identifica y corrige un
problema técnico conocido de los \gls{embeddings} tradicionales --la
anisotropía, que hace que las representaciones de distintos cursos
resulten artificialmente similares entre sí-- y constituye el precedente
más cercano encontrado al motor de \gls{matching} que \sistema{} plantea
implementar sobre perfiles de alumnos y tutores. En la misma línea, un
estudio publicado en el \emph{Journal of Applied Informatics and
Computing} \citep{paperrec2025} implementó un sistema de recomendación de
artículos científicos con sentence-transformers y similitud de coseno,
validado con una prueba de usuario de 45 personas en la que el 95,5\,\%
consideró relevantes las recomendaciones recibidas --un antecedente
metodológico útil para diseñar la validación del propio motor de
\sistema{}. Un tercer trabajo \citep{beefore2024} advierte que, hasta el
momento, el uso de sentence-transformers en sistemas de recomendación se
limita mayormente a generar información auxiliar para resolver el
problema de arranque en frío, y no a operar como mecanismo central de
recomendación --lo cual señala un espacio de aporte relativamente
inexplorado que coincide con el uso que \sistema{} le daría a esta
tecnología.

Sobre los sistemas de reputación y confianza en plataformas bilaterales,
un capítulo del libro \emph{Reengineering the Sharing Economy}
\citep{reputation} desarrolla la teoría detrás de los mecanismos de
reputación en mercados online y plantea como pregunta central de diseño
si la reputación debe ser bidireccional --como en Airbnb-- o
unidireccional --como en eBay o Amazon--, señalando que ambos
marketplaces resuelven de forma distinta el problema de la información
asimétrica entre las partes. Esta discusión respalda directamente la
decisión de diseño de \sistema{} de implementar un sistema de
calificación bidireccional entre alumnos y tutores.

Finalmente, en cuanto a la viabilidad técnica de la tutoría virtual, un
trabajo presentado en una conferencia de la ACM \citep{webrtc2025} sobre
un sistema de enseñanza remota basado en \gls{webrtc} para cursos de
inglés reportó, tras la optimización de la transmisión de audio y video
y la implementación de mecanismos de redundancia multicanal, una tasa de
pérdida de video de apenas 0,15\,\% y una latencia de pizarra de
8,2~milisegundos en un entorno de red educativa --cifras que evidencian
que es técnicamente viable sostener sesiones de videollamada estables
incluso en condiciones de conectividad variable, tal como requiere la
hipótesis central del proyecto.

\subsection{Productos y plataformas similares}

La tabla~\ref{tab:comparativa} sintetiza el relevamiento comercial de
las cinco plataformas más representativas del sector.

\begin{landscape}
    \small
    \begin{xltabular}{\linewidth}{l L L L L L}
        \caption{Comparación de plataformas de tutorías relevadas}
        \label{tab:comparativa} \\
        \toprule
        \tb{Plataforma} & \tb{Modelo de matching} & \tb{Aula virtual integrada} & \tb{Pago dentro de la plataforma} & \tb{Verificación de tutores} & \tb{Alcance / foco geográfico} \\
        \midrule
        \endfirsthead
        \toprule
        \tb{Plataforma} & \tb{Modelo de matching} & \tb{Aula virtual integrada} & \tb{Pago dentro de la plataforma} & \tb{Verificación de tutores} & \tb{Alcance / foco geográfico} \\
        \midrule
        \endhead
        \bottomrule
        \endfoot
        Preply \citep{preply} & Filtros manuales (materia, precio, disponibilidad); sin \gls{matching} semántico & Sí, parcial (Preply Classroom, hoy limitada mayormente a tutores de inglés) & Sí, por suscripción -- comisión de 100\,\% en la primera clase y de 18\,\% a 33\,\% en las siguientes & Reseñas de alumnos + video de presentación del tutor & Global; sin foco en zonas de baja densidad de oferta \\
        Wyzant \citep{wyzant} & Búsqueda por filtros & No nativa & Sí & Reseñas & Estados Unidos; más de 80.000 tutores y 4 millones de alumnos \\
        Superprof \citep{superprof} & Ninguno -- el alumno contacta manualmente a cada profesor & No & No -- el pago se coordina fuera de la plataforma & Reseñas de alumnos & Global, incluida Argentina; sin comisión para el profesor \\
        Tusclases \citep{tusclases} & Filtros (materia, precio, ubicación); sin \gls{matching} automático & No & No & Reseñas + título declarado por el profesor & España y Latinoamérica; portal líder en Argentina \\
        profe.com \citep{profecom} & Asignación manual por parte de un equipo curado, no es un \gls{marketplace} abierto & Sí, con pizarra digital & Sí & Selección propia del equipo docente de la plataforma & España; sin presencia en Argentina \\
    \end{xltabular}
\end{landscape}

Un dato relevante encontrado durante este relevamiento es que
\citep{profecom} ya ofrece a las familias un panel de seguimiento
con ``informes personalizados generados por IA con feedback detallado''
sobre el progreso del alumno --el antecedente comercial más cercano
encontrado al asistente conceptual post-sesión que \sistema{} plantea en
su alcance (sección~\ref{sec:alcance}), aunque aplicado sobre un modelo
de plantilla docente curada y no sobre un \gls{marketplace} abierto de
tutores independientes.

\textbf{Síntesis.} Ninguna de las plataformas relevadas combina
simultáneamente las cuatro piezas que propone \sistema{}: (1) un motor de
\gls{matching} basado en similitud semántica en lugar de filtros
manuales, (2) un aula virtual nativa disponible para todas las materias
y no solo para determinados perfiles de tutor, (3) un modelo de pago
protegido dentro de la plataforma, y (4) un foco explícito en ampliar la
oferta hacia zonas geográficas donde hoy es escasa o inexistente. Los
competidores de mayor escala en Argentina (Tusclases, Superprof)
replican en esencia el modelo de directorio de contactos que el
capítulo~\ref{cap:introduccion} describe como parte del problema
--búsqueda por filtros, contacto manual, sin protección de pago--
mientras que los competidores con aula virtual y pagos integrados
(Preply) no tienen foco geográfico en el país ni resuelven el
\gls{matching} de forma semántica. Ese espacio vacío es, en definitiva,
el que este proyecto busca ocupar.

\subsection{Oportunidades de diferenciación adicionales}

A partir del relevamiento de la sección anterior, se identificaron
cuatro funcionalidades que ninguna de las plataformas competidoras
ofrece y que se proponen incorporar al alcance del proyecto: el sistema
de detección de contenido inapropiado en tiempo real, la reputación con
señales implícitas, el precio dinámico por poder adquisitivo regional y
el modo de baja conectividad.

\begin{description}
    \item[Sistema de detección de contenido inapropiado en tiempo real
    (``kill-switch'' de seguridad).] Ninguna de las plataformas relevadas
    implementa moderación automática del contenido de la videollamada en
    sí misma. Existe, sin embargo, un antecedente académico directo:
    \citet{liang2013} desarrollaron SafeVchat, un sistema de detección
    automática de contenido obsceno para servicios de videollamada
    aleatoria, que combina múltiples detectores visuales mediante teoría
    de Dempster-Shafer para decidir en tiempo real si un usuario debe ser
    bloqueado. Reportaron que la combinación del sistema automático con
    moderación humana redujo el contenido ofensivo del 33,08\,\% al
    3,49\,\%. Dado que una porción significativa de los usuarios de
    \sistema{} serán menores de edad, se propone un mecanismo que analice
    localmente (\emph{edge computing}) los fotogramas mediante un modelo
    de clasificación y que ante una detección positiva corte
    automáticamente la transmisión y bloquee la cuenta preventivamente.
    El mecanismo técnico completo, incluida la verificación de edad del
    tutor mediante \gls{ocr} y DNI, se desarrolla en la
    sección~\ref{sec:seguridad}.

    \textbf{Salvedades a considerar:} ningún clasificador automático
    tiene 0\,\% de falsos positivos/negativos; es una capa adicional, no
    una garantía absoluta. El procesamiento en el dispositivo (por
    ejemplo, TensorFlow.js) evita enviar video a un servidor, facilitando
    el cumplimiento de la Ley de Protección de Datos Personales
    (Ley~25.326). El proceso de reactivación manual debe contemplar la
    notificación a un adulto responsable. Al ser un desarrollo complejo
    (visión por computadora en tiempo real), debe evaluarse su viabilidad
    dentro de los plazos del \gls{mvp}.

    \item[Reputación con fricción reducida (señales implícitas).] La
    literatura \citep{reputation} señala que las calificaciones
    explícitas (estrellas) son propensas a sesgos. Se propone combinar
    una acción mínima (reacción rápida) con señales de comportamiento
    capturadas automáticamente: tasa de sesiones completadas,
    puntualidad, y la tasa de recontratación del mismo tutor por parte
    del mismo alumno (preferencia revelada).

    \item[Precio dinámico por poder adquisitivo regional.] Se propone que
    la plataforma sugiera un rango de precio de referencia ajustado al
    poder adquisitivo de la provincia o localidad del alumno, calculado a
    partir de los datos transaccionales de la plataforma, manteniendo la
    decisión final en el tutor.

    \item[Modo de baja conectividad.] Se propone una degradación
    progresiva de calidad (video $\rightarrow$ solo audio
    $\rightarrow$ mensajería de texto) ante caídas de ancho de banda,
    sumado a la retención temporal de un buffer de video rotativo de 30
    segundos, el cual se persiste y envía como evidencia únicamente si se
    dispara el sistema de seguridad (kill-switch).
\end{description}

\subsection{Marco legal}

El diseño de la relación contractual entre la plataforma, el adulto
responsable y el menor se apoya en el sistema de capacidad progresiva
del Código Civil y Comercial de la Nación, artículos 25 y 26, que
establece que los adolescentes de entre 13 y 18 años tienen capacidad de
ejercicio creciente, pero que los contratos que los comprometan
patrimonialmente requieren la representación o el consentimiento de los
progenitores según los artículos 682 y 690, salvo los de escasa cuantía
de la vida cotidiana previstos en el artículo 684 bis, categoría en la
que un paquete de clases particulares no encuadra. De ahí surge la
decisión de que la cuenta contratante y pagadora sea siempre la del
adulto responsable.

El mismo criterio se refuerza con un antecedente impositivo: el modelo
de retención que aplican plataformas comparables en Argentina, como
Uber, Rappi, PedidosYa y Airbnb, donde AFIP obliga al agente de pago a
retener hasta un 28\,\% de lo abonado a quien no esté inscripto en el
régimen de Monotributo, funciona como antecedente directo de la
estrategia de formalización tributaria de los tutores de \sistema{},
desarrollada en el capítulo~\ref{cap:economica}. A esto se suma la
jurisprudencia laboral argentina sobre la relación de dependencia
encubierta de repartidores de plataformas de delivery, que constituye un
antecedente de riesgo legal relevante para el vínculo entre \sistema{} y
sus tutores, documentado como riesgo en el capítulo~\ref{cap:riesgos}.

\section{Bases teóricas}\label{sec:bases-teoricas}

\subsection{Variable 1: matching semántico e inteligencia artificial
aplicada a sistemas de recomendación}\label{sec:variable1}

El motor de \gls{matching} de \sistema{} se apoya conceptualmente en dos
cuerpos teóricos: los sistemas de recomendación (\emph{recommender
systems}), que buscan predecir la afinidad entre dos entidades a partir
de sus atributos; y el procesamiento de lenguaje natural (\gls{nlp})
aplicado a la representación semántica de texto mediante
\gls{embeddings}. Un \gls{embeddings} es una representación vectorial de
un texto en un espacio de alta dimensionalidad. Modelos como los
sentence-transformers generan vectores a partir de oraciones completas,
superando las búsquedas por palabras clave. Sobre este espacio,
estructuras como FAISS (Facebook AI Similarity Search) permiten
búsquedas de similitud eficientes.

Como se detalló en la sección~\ref{sec:estado-del-arte}, el trabajo de
la Universidad de Ottawa \citep{isotropy2026} validó esta técnica en la
recomendación de cursos, corrigiendo el problema de anisotropía. Este
antecedente es directamente aplicable a \sistema{}, donde los perfiles
comparten un vocabulario acotado (materias, niveles). La literatura
\citep{beefore2024} identifica que el uso de sentence-transformers como
mecanismo central sigue siendo poco explorado, confirmando que la
implementación de \sistema{} representa un aporte que excede la
aplicación estándar.

\subsection{Variable 2: acceso a educación personalizada y brecha de
aprendizaje}\label{sec:variable2}

La base teórica proviene de la investigación sobre enseñanza
individualizada de Carroll \citep{carroll1963} y Bloom
\citep{bloom1968, bloom1984}. Carroll planteó que el aprendizaje depende
de la relación entre el tiempo necesario para dominar un contenido y el
tiempo efectivamente destinado. Bloom desarrolló la teoría del
\emph{mastery learning}, sosteniendo que la mayoría puede alcanzar
niveles altos de dominio con instrucción personalizada.

El aporte central es el ``problema de los 2 sigma'' \citep{bloom1984}:
alumnos con tutoría individual rinden, en promedio, dos desviaciones
estándar por encima de alumnos en instrucción grupal convencional. La
propuesta de \sistema{} responde tecnológicamente a este problema: busca
hacer que la tutoría individual sea accesible a escala, superando la
barrera de disponibilidad geográfica. Esto es consistente con el
meta-análisis de tutoría privada \citep{sciencedirect2022} y el trabajo
sobre tutoría bajo demanda \citep{ondemand2026}, que la ubican entre las
intervenciones más efectivas.

\subsection{Relación entre las variables}

El vínculo se explica a través de la teoría económica de los mercados
bilaterales (\emph{two-sided markets}) \citep{rochet2006}. Un mercado
bilateral conecta a dos grupos (alumnos y tutores) donde la participación
de uno aumenta el valor para el otro. El mayor riesgo es el problema del
\gls{coldstart}, señalado en el capítulo~\ref{cap:introduccion}.

En un mercado geográficamente disperso (Variable~2), el problema del
\gls{coldstart} se agrava. La función del \gls{matching} semántico
(Variable~1) es el mecanismo que permite que esa masa crítica, aunque
dispersa, se encuentre eficientemente a escala nacional. Al no depender
de coincidencias exactas ni proximidad, el motor permite a \sistema{}
operar como un único mercado denso. Esta es la hipótesis técnica que
sostiene el proyecto.

\section{Marco conceptual}

\sistema{} se ubica dentro del sector \gls{edtech} y, específicamente, en
los \gls{marketplace} educativos (dimensionado en el
capítulo~\ref{cap:introduccion}).

En el plano de la infraestructura, el proyecto se apoya en \gls{webrtc},
un estándar abierto para comunicación en tiempo real en navegadores.
Sobre este, \sistema{} utiliza LiveKit, un framework open-source que
simplifica la implementación de salas escalables. Para sostener el estado
en tiempo real de esas salas, se utiliza Redis (base de datos en memoria
clave-valor), necesaria por sus exigencias de baja latencia.

En el circuito de pagos, se apoya en un modelo de \gls{escrow}
implementado a través de MercadoPago Marketplace, reteniendo los fondos
hasta la realización de la clase, reduciendo el riesgo de fraude.

Sobre la confianza, \sistema{} adopta un sistema de reputación
bidireccional, respaldado por la literatura \citep{reputation}. Para
medir la satisfacción general, se contempla el uso del \gls{nps}.
Finalmente, el proyecto se desarrolla bajo un esquema de Scrum adaptado
(capítulo~\ref{cap:metodologia}).

\section{Entrevistas a especialistas}\label{sec:entrevistas}

Se proponen dos perfiles de especialista complementarios cuyos guiones
completos se transcriben en el apéndice~\ref{apx:entrevista}:
\textbf{Perfil 1}, docente o gestor/a institucional, para indagar sobre
el rol del apoyo escolar en la prevención de la repitencia, las
diferencias urbanas-rurales en el acceso a tutorías y los recaudos
pedagógicos y de seguridad de una tutoría virtual; y \textbf{Perfil 2},
psicopedagogía / vínculo y seguridad con menores, para abordar el
impacto vincular del acompañamiento personalizado, las señales de alerta
en entornos remotos y los errores a evitar en el diseño del sistema de
reputación.

\textbf{[PENDIENTE:} una vez realizada la entrevista, se agregará acá la
síntesis de los resultados con fecha, nombre, cargo y aporte de cada
persona entrevistada. La evidencia detallada se conserva en el
apéndice~\ref{apx:entrevista}\textbf{]}
    %!TEX root = ../main.tex
% =============================================================
%  CAPÍTULO 4 — MARCO TECNOLÓGICO       (Entrega 4 · 10 % · 25/8)
% =============================================================

\chapter{Marco tecnológico}\label{cap:marco-tecnologico}

\begin{guia}[Guía de la cátedra: qué se evalúa en este capítulo]
Documenta el \emph{cómo}. No es un manual de uso: es la justificación
de cada elección tecnológica. Regla: nada ``porque sí''.
\begin{enumerate}
    \item \textbf{Arquitectura}: diagrama con componentes y conexiones,
          responsabilidad y tecnología de cada uno, flujo de datos y
          justificación del estilo arquitectónico con alternativas
          evaluadas \citep{sommerville2016}.
    \item \textbf{Prototipo o \gls{g:mvp}}: cuál es y por qué; alcance
          implementado vs.\ simulado; capturas o enlace; cómo se probó.
    \item \textbf{Stack}: por tecnología, nombre, versión, capa, por qué
          (madurez, comunidad, costo, conocimiento del equipo) y qué
          alternativas se descartaron.
    \item \textbf{Infraestructura y base de datos}: dónde se aloja cada
          componente, escalabilidad, costos (enlazados al
          capítulo~\ref{cap:economica}), motor y modelo de datos.
    \item \textbf{Diseño y \gls{ux}}: quién lo usa, mockups, decisiones de
          diseño y \textbf{cómo se validó con usuarios} (la cátedra lo marca
          como importante; cuestionario en el apéndice~\ref{apx:encuesta}).
    \item \textbf{Seguridad, calidad y despliegue}: autenticación, control
          de acceso, cifrado, plan de pruebas, repositorio, \gls{cicd} y
          monitoreo.
\end{enumerate}
\end{guia}

\completar{Desarrollo del capítulo.}

\section{\completar{Nombre de la sección}}

\completar{Texto.}

\begin{figure}[H]
    \centering
    \includegraphics[width=0.85\textwidth]{example-image-a}
    \caption{Diagrama de arquitectura del sistema (ejemplo de figura)}
    \label{fig:arquitectura}
\end{figure}

\begin{ejemplo}[Ejemplo: decisión con alternativas]
Se adoptó una arquitectura \completar{cliente-servidor en capas} porque
\completar{el volumen de usuarios previsto y el equipo de una persona no
justifican la complejidad operativa de microservicios}. Se evaluó
\completar{alternativa} y se descartó porque \completar{motivo concreto}.
\end{ejemplo}

\section{\completar{Nombre de la sección}}

\completar{Texto.}

\begin{table}[H]
    \centering
    \caption{Stack tecnológico y alternativas evaluadas (ejemplo de estructura)}
    \label{tab:stack}
    \small
    \begin{tabularx}{\textwidth}{l l l L L}
        \toprule
        \tb{Capa} & \tb{Tecnología} & \tb{Versión} & \tb{Justificación} & \tb{Alternativas descartadas} \\
        \midrule
        Frontend & \completar{framework} & \completar{x.y} & \completar{madurez, comunidad, conocimiento del equipo} & \completar{alternativa: motivo} \\
        Backend & \completar{lenguaje/framework} & \completar{x.y} & \completar{...} & \completar{...} \\
        Base de datos & \completar{motor} & \completar{x.y} & \completar{relacional por integridad referencial y consultas ad hoc} \citep{postgresqldocs} & \completar{NoSQL: sin necesidad de esquema flexible} \\
        Modelo & \completar{librería} & \completar{x.y} & \completar{...} \citep{pedregosa2011} & \completar{...} \\
        \bottomrule
    \end{tabularx}
\end{table}

\begin{figure}[H]
    \centering
    \begin{subfigure}[b]{0.48\textwidth}
        \includegraphics[width=\textwidth]{example-image-a}
        \caption{Pantalla de inicio}
    \end{subfigure}
    \hfill
    \begin{subfigure}[b]{0.48\textwidth}
        \includegraphics[width=\textwidth]{example-image-b}
        \caption{Panel principal}
    \end{subfigure}
    \caption{Mockups de pantallas (ejemplo de figura compuesta)}
    \label{fig:mockups}
\end{figure}

\begin{lstlisting}[language=Python, caption={Ejemplo de código fuente con resaltado}, label={lst:ejemplo}]
def test_calcular_exposicion():
    # exposición = probabilidad x impacto
    assert calcular_exposicion(probabilidad=4, impacto=5) == 20
\end{lstlisting}

    %!TEX root = ../main.tex
% =============================================================
%  CAPÍTULO 5 — JUSTIFICACIÓN ECONÓMICA (Entrega 5 · 15 % · 15/9)
% =============================================================

\chapter{Justificación económica}\label{cap:economica}

\begin{guia}[Guía de la cátedra: qué se evalúa en este capítulo]
Es el capítulo de mayor peso individual (15\,\%). Se evalúa como una
evaluación de proyectos de inversión:
\begin{enumerate}
    \item \textbf{Análisis de mercado y competidores}: tamaño, segmentos,
          precios de referencia.
    \item \textbf{Modelo de negocio}: B2B, B2C, B2G o interno; niveles de
          precio y fuentes de ingreso.
    \item \textbf{Inversión inicial} desglosada por rol y tiempo, más
          infraestructura.
    \item \textbf{Costos operativos} recurrentes.
    \item \textbf{Escenarios} pesimista, base y optimista, con flujo de caja
          proyectado a 3 a 5 años.
    \item \textbf{Indicadores}: \gls{van}, \gls{tir} y período de repago,
          con una tasa de descuento construida y justificada con fuentes.
    \item \textbf{Supuestos y restricciones} explícitos.
\end{enumerate}
Todo número con su origen: cotización, tarifa pública, estadística
citada o estimación propia declarada como tal.
\end{guia}

\completar{Desarrollo del capítulo.}

\section{\completar{Nombre de la sección}}

\completar{Texto.}

\begin{table}[H]
    \centering
    \caption{Inversión inicial desglosada (ejemplo de estructura)}
    \label{tab:inversion}
    \small
    \begin{tabularx}{\textwidth}{L r r r L}
        \toprule
        \tb{Concepto} & \tb{Cantidad} & \tb{Unitario (USD)} & \tb{Total (USD)} & \tb{Fuente} \\
        \midrule
        Desarrollador full stack (meses) & 6 & \completar{2.500} & \completar{15.000} & \completar{encuesta salarial citada} \\
        Científico de datos (meses) & 4 & \completar{...} & \completar{...} & \completar{...} \\
        Diseño \gls{ux} (horas) & 80 & \completar{...} & \completar{...} & \completar{...} \\
        Infraestructura inicial & 1 & \completar{...} & \completar{...} & \completar{tarifa pública del proveedor} \\
        \midrule
        \multicolumn{3}{r}{\tb{Total}} & \tb{\completar{...}} & \\
        \bottomrule
    \end{tabularx}
\end{table}

\section{\completar{Nombre de la sección}}

\completar{Texto.}

\begin{table}[H]
    \centering
    \caption{Flujo de caja proyectado por escenario, en USD (ejemplo de estructura)}
    \label{tab:flujo-caja}
    \small
    \begin{tabularx}{\textwidth}{l R R R R R R}
        \toprule
        \tb{Escenario} & \tb{Año 0} & \tb{Año 1} & \tb{Año 2} & \tb{Año 3} & \tb{Año 4} & \tb{Año 5} \\
        \midrule
        Pesimista & $-$\completar{...} & \completar{...} & \completar{...} & \completar{...} & \completar{...} & \completar{...} \\
        Base      & $-$\completar{...} & \completar{...} & \completar{...} & \completar{...} & \completar{...} & \completar{...} \\
        Optimista & $-$\completar{...} & \completar{...} & \completar{...} & \completar{...} & \completar{...} & \completar{...} \\
        \bottomrule
    \end{tabularx}
\end{table}

\begin{ejemplo}[Ejemplo: construcción de la tasa de descuento]
La tasa de descuento \gls{sym:r} se construye sumando una tasa libre de
riesgo (bono del Tesoro de EE.\,UU.\ a 10 años), el riesgo país (índice
EMBI) y una prima por riesgo de proyecto en etapa temprana. Cada
componente se cita con su fuente y fecha de consulta.
\end{ejemplo}

El \gls{sym:van} se calcula como
\begin{equation}
    \mathit{VAN} = -I_0 + \sum_{t=1}^{n} \frac{\gls{sym:fc}}{(1 + r)^{t}},
    \label{eq:van}
\end{equation}
donde $I_0$ es la inversión inicial, $FC_t$ el flujo de caja neto del
año $t$ y $r$ la tasa de descuento. La \gls{tir} es la tasa que anula el
\gls{van}, y el período de repago es el primer año en que el flujo
acumulado se vuelve positivo.

\begin{table}[H]
    \centering
    \caption{Indicadores financieros por escenario (ejemplo de estructura)}
    \label{tab:indicadores-financieros}
    \small
    \begin{tabular}{l r r r}
        \toprule
        \tb{Escenario} & \tb{VAN (USD)} & \tb{TIR} & \tb{Repago} \\
        \midrule
        Pesimista & \completar{...} & \completar{...\,\%} & \completar{n} años \\
        Base      & \completar{...} & \completar{...\,\%} & \completar{n} años \\
        Optimista & \completar{...} & \completar{...\,\%} & \completar{n} años \\
        \bottomrule
    \end{tabular}
\end{table}

    %!TEX root = ../main.tex
% =============================================================
%  CAPÍTULO 6 — ANÁLISIS DE RIESGOS     (Entrega 6 · 5 % · 29/9)
% =============================================================

\chapter{Análisis de riesgos}\label{cap:riesgos}

\begin{guia}[Guía de la cátedra: qué espera ver en este capítulo]
Formato obligatorio de cada riesgo: \textbf{``Debido a [causa], podría
ocurrir [evento], lo que provocaría [efecto]''}. Sin causa es un miedo,
no un riesgo.
\begin{enumerate}
    \item \textbf{Tabla de riesgos} con ID, descripción causa--evento--efecto,
          naturaleza, etapa, probabilidad (1--5), impacto (1--5),
          \gls{exposicion} (1--25), respuesta, responsable y contingencia;
          ordenada por exposición descendente. Con los 6 a 8 primeros bien
          tratados alcanza; el resto se monitorea.
    \item Al menos un riesgo de \textbf{cada naturaleza} (alcance, técnicos,
          recursos, terceros, datos y cumplimiento) y de \textbf{cada etapa}
          (antes de arrancar, durante, pasaje a producción, operación).
    \item Al menos un riesgo \textbf{externo} (``si no tienen ninguno, no
          miraron para afuera'') y al menos uno con respuesta
          \textbf{aceptar} bien justificado (``aceptar se escribe; ignorar
          se paga'').
    \item \textbf{Supuestos y restricciones} separados de los riesgos. Cada
          supuesto que falla se convierte en riesgo.
    \item Que sean propios: ``un riesgo copiado de internet se nota a tres
          metros''.
\end{enumerate}
Los siete que matan proyectos: alcance que nunca cierra (el número uno
de la materia), pérdida de la fuente de datos, equipo que se desarma,
tecnología que no da, inviabilidad de costo, barrera legal o ética,
nadie a cargo. Y la pregunta de jurado: ¿quién lo mantiene cuando
ustedes se reciben? Marco de referencia: \citet{pmi2021}.
\end{guia}

\section{Enfoque y criterios}

Un riesgo es un evento incierto que, de ocurrir, afecta al menos a uno
de los objetivos del proyecto; a diferencia de un problema, todavía no
ocurrió y admite decisiones anticipadas \citep{pmi2021}. Cada riesgo de
este capítulo se formula como ``debido a [causa], podría ocurrir
[evento], lo que provocaría [efecto]'', se clasifica por su naturaleza
(alcance, técnico, recursos, terceros, datos o cumplimiento), por su
origen (interno, bajo control del proyecto, o externo) y por la etapa en
que puede materializarse (antes de arrancar, durante el desarrollo, en
el pasaje a producción o en la operación).

La \gls{exposicion} se calcula como el producto de la probabilidad y el
impacto, ambos en una escala de 1 a 5
(tabla~\ref{tab:exposicion}). Para este proyecto, las escalas se
interpretan de la siguiente manera:

\begin{itemize}
    \item \textbf{Probabilidad.} 1: sin indicios de que ocurra; 3: hay
          evidencia de que la causa existe; 5: la causa ya está presente y
          nada la contiene.
    \item \textbf{Impacto.} 1: demora de días sin efecto sobre los
          objetivos; 3: afecta un objetivo específico o reduce el
          \gls{van} sin volverlo negativo; 5: impide el lanzamiento,
          vuelve negativo el \gls{van} del escenario base o expone a
          menores a un daño.
\end{itemize}

Las respuestas posibles son cuatro: \textbf{evitar} (eliminar la causa,
cambiando el alcance o la actividad), \textbf{mitigar} (reducir la
probabilidad o el impacto), \textbf{transferir} (trasladar el riesgo a un
tercero mejor preparado) y \textbf{aceptar} (convivir con él,
documentado y con un plan de contingencia).

\section{Identificación y evaluación}

La tabla~\ref{tab:matriz-riesgos} presenta los quince riesgos
identificados, ordenados por exposición descendente. Las fuentes de
identificación fueron la encuesta de validación
(apéndice~\ref{apx:encuesta}), el modelo financiero
(capítulo~\ref{cap:economica}), el estado del repositorio
(anexo~\ref{anx:repositorios}) y las condiciones de los proveedores
relevadas en el capítulo~\ref{cap:marco-tecnologico}.

\begin{table}[H]
    \centering
    \caption{Ubicación de los riesgos en la matriz de exposición (probabilidad $\times$ impacto)}
    \label{tab:exposicion}
    \small
    \setlength{\tabcolsep}{4pt}
    \begin{tabular}{l c c c c c}
        \toprule
        & \tb{1 Insignif.} & \tb{2 Menor} & \tb{3 Moderado} & \tb{4 Importante} & \tb{5 Severo} \\
        \midrule
        \tb{5 Muy probable}   & \riesgoBajo{}  & \riesgoMedio{} & \riesgoAlto{} & \riesgoAlto{} & \riesgoAlto{} \\
        \tb{4 Probable}       & \riesgoBajo{}  & \riesgoMedio{} & \riesgoAlto{R-06, R-08} & \riesgoAlto{R-02, R-03} & \riesgoAlto{R-01} \\
        \tb{3 Posible}        & \riesgoBajo{}  & \riesgoMedio{R-13, R-14} & \riesgoMedio{R-11} & \riesgoAlto{R-07} & \riesgoAlto{R-04, R-05} \\
        \tb{2 Poco probable}  & \riesgoBajo{}  & \riesgoBajo{}  & \riesgoMedio{R-15} & \riesgoMedio{R-12} & \riesgoMedio{R-09, R-10} \\
        \tb{1 Muy improbable} & \riesgoBajo{}  & \riesgoBajo{}  & \riesgoBajo{} & \riesgoBajo{} & \riesgoBajo{} \\
        \bottomrule
    \end{tabular}
\end{table}

% La matriz va apaisada (landscape) porque tiene muchas columnas.
\begin{landscape}
\small
\begin{xltabular}{\linewidth}{l >{\raggedright\arraybackslash\hsize=1.3\hsize}X >{\raggedright\arraybackslash}p{2.3cm} >{\raggedright\arraybackslash}p{1.7cm} c c c >{\raggedright\arraybackslash\hsize=0.85\hsize}X >{\raggedright\arraybackslash}p{1.5cm} >{\raggedright\arraybackslash\hsize=0.85\hsize}X}
    \caption{Matriz de riesgos ordenada por exposición}
    \label{tab:matriz-riesgos} \\
    \toprule
    \tb{ID} & \tb{Descripción (causa, evento, efecto)} & \tb{Naturaleza (origen)} & \tb{Etapa} & \tb{P} & \tb{I} & \tb{E} & \tb{Respuesta} & \tb{Resp.} & \tb{Contingencia} \\
    \midrule
    \endfirsthead
    \toprule
    \tb{ID} & \tb{Descripción (causa, evento, efecto)} & \tb{Naturaleza (origen)} & \tb{Etapa} & \tb{P} & \tb{I} & \tb{E} & \tb{Respuesta} & \tb{Resp.} & \tb{Contingencia} \\
    \midrule
    \endhead
    \bottomrule
    \endfoot
    R-01 & Debido a que el resumen automático envía el audio de la sesión directamente al \gls{llm} de un proveedor extranjero, y la anonimización prevista opera sobre texto, podría ocurrir que la voz y los datos personales de menores se transfieran sin anonimizar, lo que provocaría un incumplimiento de la Ley 25.326 y la pérdida de confianza de las familias. & Cumplimiento (interno) & Durante & 4 & 5 & \riesgoAlto{20} & Evitar: en el \gls{mvp} el adicional de resumen no se ofrece en sesiones con menores; para adultos, consentimiento explícito previo & Autor & Si se detecta un envío indebido: desactivar el módulo de resumen, solicitar la eliminación al proveedor y notificar a los titulares \\
    R-02 & Debido a que solo el 35,2\,\% de los encuestados prefiere pagar dentro de una plataforma antes que por transferencia, podría ocurrir que alumnos y tutores acuerden las clases por fuera después del primer contacto, lo que provocaría la pérdida de la comisión, principal fuente de ingreso. & Terceros: usuarios (externo) & Operación & 4 & 4 & \riesgoAlto{16} & Mitigar: \gls{escrow}, verificación, reputación y resumen solo existen dentro de la plataforma; datos de contacto ocultos hasta la primera reserva & Autor & Si la recompra dentro de la plataforma cae por debajo del 50\,\%: comisión decreciente por fidelidad del tutor \\
    R-03 & Debido a que el capítulo~\ref{cap:conclusiones} vence el 20 de octubre y el piloto se extiende del 27 de octubre al 27 de noviembre, con 36 tareas pendientes a cargo de un único desarrollador, podría ocurrir que la hipótesis deba responderse sin datos del piloto, lo que provocaría conclusiones sin evidencia suficiente. & Alcance (interno) & Durante & 4 & 4 & \riesgoAlto{16} & Mitigar: alcance del \gls{mvp} congelado; pruebas con usuarios reducidas antes del 20 de octubre & Autor & Presentar la respuesta a la hipótesis como preliminar y completarla con el piloto para la entrega final \\
    R-04 & Debido a que el escenario pesimista no cubre los costos fijos en tres de los cinco años, podría ocurrir que la demanda no supere las 5.100 sesiones anuales, lo que provocaría que la inversión no se recupere (\gls{van} de \usd{-5.468}). & Recursos: mercado (externo) & Operación & 3 & 5 & \riesgoAlto{15} & Mitigar: piso de \usd{4} por hora, presupuesto de adquisición y acuerdos institucionales & Autor & Si a los 12 meses el volumen anualizado es menor a 5.100 sesiones: pasar a clases grupales o al segmento de ingreso universitario; si persiste, cierre ordenado \\
    R-05 & Debido a que el clasificador del kill-switch no está integrado en el cliente (T-M3-06), podría ocurrir que el piloto incluya sesiones con menores sin la protección prevista, lo que provocaría la exposición de menores a contenido inapropiado. & Técnico (interno) & Pasaje a producción & 3 & 5 & \riesgoAlto{15} & Evitar: T-M3-06 es criterio de entrada para cualquier sesión con un menor & Autor & Piloto limitado a alumnos mayores de edad \\
    R-06 & Debido a que los proveedores de \gls{llm} discontinúan modelos y modifican precios (Gemini 2.0 Flash fue dado de baja en junio de 2026 y Gemini 3.6 Flash duplica su precio en enero de 2027), podría ocurrir que el modelo elegido deje de estar disponible o se encarezca, lo que provocaría la interrupción del resumen o la pérdida de su margen. & Terceros (externo) & Operación & 4 & 3 & \riesgoAlto{12} & Mitigar: el backend accede al proveedor a través de un puerto intercambiable (\texttt{ResumenProveedor}) & Autor & Cambiar de proveedor y ajustar el precio del adicional; si no es viable, suspenderlo sin afectar el núcleo \\
    R-07 & Debido a que solo 9 tutores respondieron la encuesta y el \gls{marketplace} es bilateral, podría ocurrir que no haya suficientes tutores verificados al iniciar el piloto, lo que provocaría búsquedas sin candidatos e impediría medir la aceptación del \gls{matching} (I-01). & Recursos: oferta (externo) & Antes de arrancar & 3 & 4 & \riesgoAlto{12} & Mitigar: reclutar tutores antes del piloto en la comunidad universitaria & Autor & Acotar el piloto a una o dos materias con oferta suficiente \\
    R-08 & Debido a que el 90,9\,\% de las respuestas de la encuesta proviene del AMBA, podría ocurrir que la demanda y el precio del interior difieran de los relevados, lo que provocaría que la hipótesis de reducción de la brecha geográfica quede sin evidencia. & Datos (interno) & Antes de arrancar & 4 & 3 & \riesgoAlto{12} & Mitigar: reencuesta dirigida al interior y entrevista a docentes de zonas rurales & Autor & Declarar la limitación y acotar el alcance de la conclusión al AMBA \\
    R-09 & Debido a que el desarrollo y la operación dependen de una sola persona, podría ocurrir una enfermedad, una sobrecarga o un abandono, lo que provocaría la detención del desarrollo o de la operación. & Recursos (interno) & Operación & 2 & 5 & \riesgoMedio{10} & Aceptar: el modelo económico no admite un equipo rentado (capítulo~\ref{cap:economica}); documentación y pruebas reducen la dependencia & Autor & Prórroga ante la cátedra; en operación, pausa sin pérdida de datos y traspaso a partir del repositorio documentado \\
    R-10 & Debido a que la verificación de antecedentes depende de la revisión manual del \gls{cap}, un documento que presenta el propio tutor, podría ocurrir que un certificado adulterado o evaluado con un criterio incorrecto habilite a un tutor con antecedentes para dar clases a menores, lo que provocaría un daño grave al menor y responsabilidad legal para la plataforma. & Cumplimiento (interno) & Operación & 2 & 5 & \riesgoMedio{10} & Mitigar: \gls{cap} vigente obligatorio para dictar clases a menores, con verificación de la firma digital del Registro Nacional de Reincidencia y criterio de rechazo definido con asesoría legal & Autor y asesoría legal & Suspensión inmediata, preservación de la evidencia y denuncia ante la autoridad \\
    R-11 & Debido a que el cobro con MercadoPago solo se probó con usuarios de prueba, dado que el entorno sandbox presenta fallas conocidas en las notificaciones, podría ocurrir que en producción los webhooks lleguen tarde, duplicados o no lleguen, lo que provocaría reservas confirmadas sin cobro o cobros sin reserva. & Técnico: terceros (externo) & Pasaje a producción & 3 & 3 & \riesgoMedio{9} & Mitigar: procesamiento idempotente de las notificaciones y conciliación periódica contra la \gls{api} & Autor & Modo bypass (ADR-M5-01) y conciliación manual por el rol de soporte financiero \\
    R-12 & Debido a la jurisprudencia laboral sobre repartidores de plataformas de delivery, podría ocurrir que un tutor reclame una relación de dependencia con la plataforma, lo que provocaría costos laborales no previstos. & Cumplimiento (externo) & Operación & 2 & 4 & \riesgoMedio{8} & Transferir: términos y condiciones redactados por asesoría legal; el tutor fija precio y horarios & Asesoría legal & Defensa legal con los registros de autonomía del tutor \\
    R-13 & Debido a que la infraestructura se contrata en dólares a proveedores que pueden ajustar sus tarifas (Hetzner lo hizo dos veces en 2026), podría ocurrir un aumento de los costos fijos, lo que provocaría una reducción del \gls{van}. & Terceros (externo) & Operación & 3 & 2 & \riesgoMedio{6} & Mitigar: despliegue en contenedores portables entre proveedores & Autor & Migrar de proveedor con la misma configuración de Coolify \\
    R-14 & Debido a que el esfuerzo de las 36 tareas pendientes se estima con un promedio histórico de 0,76 horas por tarea, podría ocurrir que las tareas más complejas demanden más tiempo, lo que provocaría un aumento de la inversión y un retraso del piloto. & Recursos (interno) & Durante & 3 & 2 & \riesgoMedio{6} & Aceptar dentro del margen: el \gls{van} base tolera hasta 64 horas adicionales & Autor & Cubrir el desvío con la contingencia del 10\,\% y postergar las tareas no críticas (T-M4-12 a T-M4-15) \\
    R-15 & Debido a que el backend se migra de la computadora del desarrollador a un \gls{vps} antes del piloto, podría ocurrir una falla en el despliegue o en las migraciones de esquema, lo que provocaría una interrupción del servicio al inicio del piloto. & Técnico (interno) & Pasaje a producción & 2 & 3 & \riesgoMedio{6} & Mitigar: migraciones versionadas con Flyway y despliegue ensayado antes del piloto & Autor & Volver a la imagen Docker anterior y restaurar el último respaldo \\
\end{xltabular}
\end{landscape}

\section{Tratamiento de los riesgos principales}

Los diez primeros riesgos concentran la exposición del proyecto. Su
tratamiento se detalla a continuación; el resto se monitorea.

\begin{itemize}
    \item \textbf{R-01, privacidad del audio de menores (20).} La causa
          está presente en el diseño: el Plan M6 del repositorio prevé
          anonimizar el texto antes de enviarlo al \gls{llm}, pero el
          proveedor elegido recibe el audio directamente
          (capítulo~\ref{cap:economica}). Por el impacto, se opta por
          evitar el riesgo en lugar de mitigarlo: el adicional de resumen
          queda fuera de las sesiones con menores hasta que exista un
          paso de transcripción previo a la anonimización. La decisión
          reduce la adopción esperada del adicional, supuesto que se
          revisa en la sección~\ref{sec:supuestos}.
    \item \textbf{R-02, desintermediación (16).} Es el riesgo con mejor
          respaldo empírico: el 34,1\,\% de los encuestados se mostró en
          desacuerdo con pagar dentro de una plataforma
          (apéndice~\ref{apx:encuesta}). La respuesta apunta a que el
          valor diferencial solo exista dentro de \sistema{}: las
          funciones más valoradas en la encuesta ---profesores
          verificados (58\,\%), resumen automático (56,8\,\%) y pago
          seguro (48,9\,\%)--- dependen de que la transacción ocurra en
          la plataforma.
    \item \textbf{R-03, alcance y cronograma (16).} Corresponde al
          ``alcance que nunca cierra'' señalado por la cátedra. El
          alcance del \gls{mvp} se congela en las 36 tareas pendientes, y
          cualquier funcionalidad nueva pasa al trabajo futuro.
    \item \textbf{R-04, inviabilidad económica (15).} Surge del escenario
          pesimista del capítulo~\ref{cap:economica}. La contingencia
          define un disparador medible y una decisión: si a los doce meses
          el volumen anualizado no alcanza el punto de equilibrio
          operativo, se revisa el modelo de negocio y, si no hay
          alternativa viable, se cierra la operación de forma ordenada.
          Cancelar a tiempo un proyecto inviable es una decisión de
          gestión, no un fracaso.
    \item \textbf{R-05, kill-switch sin integrar (15).} El backend del
          kill-switch está completo, pero el clasificador que corre en el
          navegador no. La respuesta es evitar: sin esa integración no se
          habilitan sesiones con menores.
    \item \textbf{R-06, proveedor de \gls{llm} (12).} La probabilidad es
          alta porque ya ocurrió durante el desarrollo: el candidato
          original fue dado de baja. El impacto es moderado porque el
          resumen es un adicional y el backend lo aísla detrás de un
          puerto intercambiable.
    \item \textbf{R-07, oferta de tutores (12).} Es la manifestación del
          problema de arranque en frío (\gls{coldstart}) descrito en el
          capítulo~\ref{cap:marco-teorico}.
    \item \textbf{R-08, sesgo geográfico de la encuesta (12).} Afecta
          directamente a la hipótesis central, que pone el foco en las
          zonas con escasa oferta de tutores.
    \item \textbf{R-09, dependencia de una persona (10).} Se acepta
          porque no existe una respuesta alternativa compatible con el
          proyecto: el análisis económico mostró que un equipo rentado
          torna negativo el \gls{van} en los tres escenarios. Aceptarlo no
          equivale a ignorarlo: la documentación de diseño (especificaciones,
          planes y decisiones de arquitectura) y 455 casos de prueba
          automatizados permiten que otra persona retome el sistema.
    \item \textbf{R-10, antecedentes de los tutores de menores (10).}
          Una primera versión del sistema exigía el \gls{cap} a todos los
          tutores y se retiró (ADR-M1-02) por dos motivos: la fricción en
          el registro de tutores, el recurso más escaso de la plataforma
          (R-07), y la necesidad de decidir la descalificación de un tutor
          por sus antecedentes sin asesoría legal. La verificación se
          reincorpora acotada: el \gls{cap} vigente es obligatorio solo
          para dictar clases a menores. Así, los tutores que enseñan a
          adultos se registran sin fricción adicional, y el criterio de
          rechazo se define con la asesoría legal ya presupuestada
          (capítulo~\ref{cap:economica}). El riesgo residual es la
          validez del documento y su evaluación manual.
\end{itemize}

\section{Riesgos que cancelan proyectos}

La tabla~\ref{tab:riesgos-fatales} vincula los siete riesgos que la
cátedra señala como causas de cancelación con los riesgos identificados.

\begin{table}[H]
    \centering
    \caption{Riesgos fatales y su tratamiento en \sistema{}}
    \label{tab:riesgos-fatales}
    \small
    \begin{tabularx}{\textwidth}{l L}
        \toprule
        \tb{Riesgo fatal} & \tb{Situación en \sistema{}} \\
        \midrule
        Alcance que no cierra     & R-03: alcance congelado en las 36 tareas pendientes \\
        Pérdida del proveedor de datos & R-06 y R-11: los servicios de terceros se aíslan detrás de interfaces intercambiables y existe un modo de operación sin cobro real \\
        Equipo que se desarma     & R-09: aceptado, con documentación y pruebas que permiten el traspaso \\
        Tecnología inadecuada     & R-05: la tecnología existe y el spike fue exitoso (ADR-M3-01); falta su integración \\
        Inviabilidad económica    & R-04: viable en los escenarios base y optimista, con disparador de revisión \\
        Barrera ética o legal     & R-01 se evita; R-10 se mitiga con el \gls{cap} obligatorio para dictar clases a menores; R-12 se transfiere \\
        Falta de un responsable   & El autor es responsable de todos los riesgos internos; los de cumplimiento se comparten con la asesoría legal \\
        \bottomrule
    \end{tabularx}
\end{table}

\section{Oportunidades}

El análisis también identificó dos riesgos positivos, que se gestionan
buscando que ocurran:

\begin{itemize}
    \item \textbf{O-01.} Debido a que el 56,8\,\% de los encuestados
          valoró el resumen automático, podría ocurrir que la adopción
          del adicional supere el 40\,\% supuesto, lo que provocaría un
          margen por sesión mayor al proyectado. Respuesta: explotar,
          destacando el adicional en el flujo de reserva.
    \item \textbf{O-02.} Debido a que algunas jurisdicciones ya recurren a
          escuelas mediadas por tecnología con docentes en las ciudades
          \citep{olea}, podría ocurrir que una institución educativa o un
          municipio contrate paquetes de sesiones, lo que provocaría un
          volumen agregado superior al del escenario optimista. Respuesta:
          mejorar la probabilidad, presentando el piloto a instituciones
          del interior.
\end{itemize}

\section{Supuestos y restricciones}\label{sec:supuestos}

Los supuestos son condiciones que se dan por ciertas para que el
proyecto sea válido; si alguno falla, se convierte en un riesgo. Las
restricciones son condiciones impuestas que limitan el tiempo, el costo o
el alcance.

\begin{table}[H]
    \centering
    \caption{Supuestos del proyecto y riesgo asociado si fallan}
    \label{tab:supuestos}
    \small
    \begin{tabularx}{\textwidth}{l L l}
        \toprule
        \tb{ID} & \tb{Supuesto} & \tb{Si falla} \\
        \midrule
        S-01 & Los alumnos activos por año se ubican en la banda del escenario base (200 a 290) & R-04 \\
        S-02 & Cada alumno activo toma tres clases por mes & R-04 \\
        S-03 & La mediana de disposición a pagar (\usd{4,56} por hora) es representativa fuera del AMBA & R-08 \\
        S-04 & El 40\,\% de las sesiones contrata el adicional de resumen; con la exclusión de las sesiones con menores (R-01), el supuesto debe verificarse en el piloto & R-04 \\
        S-05 & Las tarifas de los proveedores se mantienen constantes en dólares durante el horizonte & R-13 \\
        S-06 & El plan gratuito de Upstash alcanza para el volumen de chat y \emph{rate limiting} & R-13 \\
        S-07 & Las tareas pendientes demandan, en promedio, el esfuerzo observado en el repositorio & R-14 \\
        S-08 & La mayoría de los pagos se realiza dentro de la plataforma & R-02 \\
        \bottomrule
    \end{tabularx}
\end{table}

\begin{table}[H]
    \centering
    \caption{Restricciones del proyecto}
    \label{tab:restricciones}
    \small
    \begin{tabularx}{\textwidth}{l L}
        \toprule
        \tb{Tipo} & \tb{Restricción} \\
        \midrule
        Tiempo    & Presentación final el 28 de noviembre de 2026, con entregas parciales por capítulo \\
        Recursos  & Un único desarrollador, con una dedicación de 15 a 20 horas semanales \\
        Costo     & Entre 0 y 100 USD mensuales durante el desarrollo y el piloto \\
        Técnica   & Los organismos estatales (RENAPER, Registro Nacional de Reincidencia) no ofrecen \gls{api} públicas \\
        Legal     & Ley 25.326 de protección de datos personales; tutores exclusivamente mayores de edad \\
        Alcance   & \gls{pwa} sin aplicaciones móviles nativas \\
        \bottomrule
    \end{tabularx}
\end{table}

\section{Seguimiento y continuidad}

Los riesgos se revisan en cada Sprint Review (capítulo~\ref{cap:metodologia}):
se actualizan la probabilidad y el impacto, se incorporan los riesgos
nuevos y se registran los que se materializaron como problemas. Durante
la operación, los disparadores de las contingencias ---el volumen
anualizado de sesiones para R-04 y la tasa de recompra dentro de la
plataforma para R-02--- se miden mensualmente.

En cuanto a la pregunta de quién mantiene el sistema una vez finalizado
el proyecto académico: el modelo económico prevé que el equipo fundador
asuma el mantenimiento correctivo y evolutivo durante los cinco años
evaluados, sin personal rentado (capítulo~\ref{cap:economica}). Si el
proyecto no continuara, el cierre se haría de forma ordenada: aviso a
los usuarios, reembolso de las reservas pendientes a través de
MercadoPago, y exportación o eliminación de los datos personales
conforme a la Ley 25.326.

    %!TEX root = ../main.tex
% =============================================================
%  CAPÍTULO 7 — PRESENTACIÓN DE RESULTADOS
%                                       (Entrega 7 · 5 % · 13/10)
% =============================================================

\chapter{Presentación de resultados}\label{cap:resultados}

\begin{guia}[Guía de la cátedra: qué se evalúa en este capítulo]
\begin{enumerate}
    \item \textbf{Requerimientos formales}: \gls{rf} y \gls{rnf} con ID,
          descripción, prioridad y criterio de aceptación medible,
          trazados a los indicadores del capítulo~\ref{cap:metodologia}.
    \item \textbf{Evidencia de implementación}: capturas reales del sistema
          funcionando, enlaces al repositorio (anexo~\ref{anx:repositorios})
          y al despliegue.
    \item \textbf{Resultados cuantitativos honestos}: diseño experimental,
          métricas, matrices de confusión o su equivalente en el dominio,
          incluyendo lo que salió mal y su análisis. Reportar también lo
          malo y argumentarlo vale más que esconderlo. Hiperparámetros y
          versiones en el apéndice~\ref{apx:tecnicos}.
    \item \textbf{Validación con usuarios y expertos}: idealmente
          triangulada (métricas técnicas, encuesta, entrevista) y
          traducida a decisiones de diseño concretas. Instrumentos en los
          apéndices \ref{apx:encuesta} y \ref{apx:entrevista}.
\end{enumerate}
El sistema tiene que existir y andar: la entrega final incluye el
proyecto funcional.
\end{guia}

\completar{Desarrollo del capítulo.}

\section{\completar{Nombre de la sección}}

\completar{Texto.}

\begin{table}[H]
    \centering
    \caption{Requerimientos funcionales (ejemplo de estructura)}
    \label{tab:rf}
    \small
    \begin{tabularx}{\textwidth}{l L L l L}
        \toprule
        \tb{ID} & \tb{Nombre} & \tb{Descripción} & \tb{Prioridad} & \tb{Criterio de aceptación} \\
        \midrule
        RF-01 & Autenticación & El sistema permite registrarse e iniciar sesión con credenciales propias & Alta & Un usuario nuevo completa el registro y accede en menos de 2 minutos \\
        RF-02 & \completar{...} & \completar{...} & \completar{...} & \completar{...} \\
        \bottomrule
    \end{tabularx}
\end{table}

\begin{table}[H]
    \centering
    \caption{Requerimientos no funcionales (ejemplo de estructura)}
    \label{tab:rnf}
    \small
    \begin{tabularx}{\textwidth}{l L L l L}
        \toprule
        \tb{ID} & \tb{Atributo} & \tb{Descripción} & \tb{Prioridad} & \tb{Umbral medible} \\
        \midrule
        RNF-01 & Rendimiento & Tiempo de respuesta de las consultas principales & Alta & Latencia p95 menor a 5 s (indicador I-02) \\
        RNF-02 & Calidad del modelo & Capacidad predictiva sobre datos no vistos & Alta & $F_1 \geq 0{,}70$ (indicador I-01) \\
        RNF-03 & \completar{...} & \completar{...} & \completar{...} & \completar{...} \\
        \bottomrule
    \end{tabularx}
\end{table}

\section{\completar{Nombre de la sección}}

\completar{Texto.}

\begin{figure}[H]
    \centering
    \includegraphics[width=0.85\textwidth]{example-image-c}
    \caption{Captura del sistema en funcionamiento (ejemplo de figura)}
    \label{fig:captura}
\end{figure}

\begin{table}[H]
    \centering
    \caption{Resultados sobre el conjunto de prueba (ejemplo de estructura)}
    \label{tab:metricas}
    \small
    \begin{tabular}{l r r r r}
        \toprule
        \tb{Modelo} & \tb{Precisión} & \tb{Exhaustividad} & \tb{$F_1$} & \tb{Tiempo (s)} \\
        \midrule
        Línea de base & \completar{...} & \completar{...} & \completar{...} & \completar{...} \\
        Modelo propuesto & \completar{...} & \completar{...} & \completar{...} & \completar{...} \\
        \bottomrule
    \end{tabular}
\end{table}

\begin{table}[H]
    \centering
    \caption{Traducción de la retroalimentación en decisiones de diseño (ejemplo de estructura)}
    \label{tab:feedback-decisiones}
    \small
    \begin{tabularx}{\textwidth}{l L L l}
        \toprule
        \tb{Fuente} & \tb{Hallazgo} & \tb{Decisión} & \tb{Estado} \\
        \midrule
        Encuesta (n=\completar{..}) & \completar{El 60\,\% no encontró la función X} & \completar{Mover X al menú principal} & Implementado \\
        Experto 1 & \completar{...} & \completar{...} & Planificado \\
        \bottomrule
    \end{tabularx}
\end{table}

    %!TEX root = ../main.tex
% =============================================================
%  CAPÍTULO 8 — IMPLICANCIAS, CONCLUSIONES Y RECOMENDACIONES
%                                       (Entrega 8 · 5 % · 20/10)
% =============================================================

\chapter{Implicancias, conclusiones y recomendaciones}\label{cap:conclusiones}

\begin{guia}[Guía de la cátedra: qué se evalúa en este capítulo]
\begin{enumerate}
    \item \textbf{Respuesta explícita a la hipótesis} del
          capítulo~\ref{cap:introduccion}, con matices y contraejemplos
          si corresponde (``afirmativa, con matices'' es una respuesta
          válida; ``se cumplieron todos los objetivos'' sin evidencia no).
    \item \textbf{Implicancias}: técnicas, sociales y económicas.
    \item \textbf{Limitaciones honestas}, por categoría (técnicas, de datos,
          de validación, de mercado), sin maquillaje. Declarar como
          preliminar lo que es preliminar.
    \item \textbf{Recomendaciones} accionables y \textbf{trabajo futuro}.
    \item \textbf{Contribución al conocimiento}: qué sabe ahora la disciplina
          que antes no sabía.
\end{enumerate}
Cierra el arco narrativo: cada objetivo específico se retoma y se
indica en qué medida se alcanzó, con referencia a la sección que lo
demuestra.
\end{guia}

\completar{Desarrollo del capítulo.}

\section{\completar{Nombre de la sección}}

\begin{ejemplo}[Ejemplo: forma de la respuesta a la hipótesis]
La hipótesis planteada en el capítulo~\ref{cap:introduccion} se responde
de manera \completar{afirmativa, con matices}: el sistema alcanzó
\completar{métrica} sobre el umbral definido en el
capítulo~\ref{cap:metodologia}, aunque \completar{limitación concreta}.
\end{ejemplo}

\completar{Texto.}

\begin{table}[H]
    \centering
    \caption{Grado de cumplimiento de los objetivos específicos (ejemplo de estructura)}
    \label{tab:cumplimiento}
    \small
    \begin{tabularx}{\textwidth}{l L l L}
        \toprule
        \tb{Objetivo} & \tb{Evidencia} & \tb{Grado} & \tb{Observación} \\
        \midrule
        OE-1 & Capítulo~\ref{cap:resultados} & Alcanzado & \completar{...} \\
        OE-2 & \completar{...} & Parcial & \completar{...} \\
        \bottomrule
    \end{tabularx}
\end{table}


    % bibliografía
    \backmatter

    % apéndices: material producido por el autor (A, B, C...)
    \begin{appendices}
        %!TEX root = ../main.tex
% APÉNDICE A — Encuesta de validación

\chapter{Encuesta de validación}\label{apx:encuesta}

\begin{guia}
Los apéndices contienen material producido por el autor que respalda
el cuerpo pero lo interrumpiría. Acá va la encuesta completa: para qué
se hizo, cómo se diseñó y distribuyó, a cuántas personas llegó y el
cuestionario tal como lo vieron quienes respondieron. Los resultados
se analizan en el capítulo~\ref{cap:resultados}.
\end{guia}

\section{Introducción}

\completar{Objetivo de la encuesta y qué se quería validar.}

\section{Metodología}

\completar{Diseño del cuestionario, población, canal de distribución,
período de aplicación y cantidad de respuestas.}

\section{Presentación de la encuesta}

\completar{Texto introductorio que vieron las personas participantes
(propósito, anonimato, duración estimada).}

\section{Cuestionario}

\begin{enumerate}
    \item \completar{Pregunta 1 (tipo: opción única).}
    \item \completar{Pregunta 2 (tipo: escala 1--5).}
    \item \completar{Pregunta 3 (tipo: abierta).}
\end{enumerate}

        %!TEX root = ../main.tex
% APÉNDICE B — Entrevistas a especialistas: guiones
% Los guiones completos se referencian desde el cap. 3 (3.4).

\chapter{Entrevistas a especialistas}\label{apx:entrevista}

\begin{guia}
Apéndice por entrevista: contexto (fecha, medio, duración, objetivo),
datos de la persona entrevistada (formación, cargo, trayectoria, con su
consentimiento) y la transcripción organizada por temas. Las
conclusiones van en el cuerpo (capítulos \ref{cap:marco-teorico} y
\ref{cap:resultados}); acá va la evidencia.
\end{guia}

Como se detalla en el capítulo~\ref{cap:marco-teorico}
(sección~\ref{sec:entrevistas}), se definieron dos perfiles de
especialista complementarios: un docente o gestor/a institucional y un/a
profesional del campo psicopedagógico o del vínculo y la seguridad con
menores. En cada caso se documenta el contexto, la ficha de la persona
entrevistada y la transcripción por ejes temáticos.

\textbf{Estado: [PENDIENTE]} -- el cronograma del
capítulo~\ref{cap:metodologia} prevé su realización durante el Sprint 3.
Hasta entonces los guiones aquí transcritos constituyen la evidencia del
diseño metodológico; los \completar{completar} marcan los datos que se
completarán con la realización efectiva de cada entrevista.

\section{Perfil 1: Docente o gestor/a institucional}

\subsection{Contexto de la entrevista}

\completar{Fecha, modalidad (presencial/remota), medio y duración.}

\subsection{Datos de la persona entrevistada}

\completar{Formación, cargo actual, experiencia relevante.}

\subsection{Guion de entrevista}

\emph{Eje 1 -- apoyo escolar y repitencia.}

\begin{enumerate}
    \item ¿Qué rol juega el apoyo escolar personalizado en la prevención
          de la repitencia en los niveles primario y secundario?
    \item ¿Qué diferencias observa entre el acceso a tutorías en zonas
          urbanas y rurales del país?
\end{enumerate}

\emph{Eje 2 -- tutoría virtual pedagógica y de seguridad.}

\begin{enumerate}
    \item Desde lo pedagógico, ¿qué recaudos considera necesarios para que
          una tutoría virtual 1 a 1 sea efectiva con un menor?
    \item Desde la seguridad, ¿qué condiciones mínimas debería garantizar
          una plataforma que conecta adultos desconocidos con menores en
          un entorno remoto?
    \item ¿Cómo observaría la articulación de esta plataforma con las
          instituciones educativas (escuelas, docentes, equipos de
          orientación)?
\end{enumerate}

\subsection{Transcripción de la entrevista}

\completar{Transcripción organizada por ejes temáticos.}

\section{Perfil 2: Psicopedagogía / vínculo y seguridad con menores}

\subsection{Contexto de la entrevista}

\completar{Fecha, modalidad (presencial/remota), medio y duración.}

\subsection{Datos de la persona entrevistada}

\completar{Formación, cargo actual, experiencia relevante.}

\subsection{Guion de entrevista}

\emph{Eje 1 -- vínculo y acompañamiento.}

\begin{enumerate}
    \item ¿Qué impacto vincular observa en un acompañamiento personalizado
          1 a 1 para niños, niñas y adolescentes?
    \item ¿Qué señales de alerta de un menor conviene monitorear en un
          entorno remoto de interacción con un adulto no conocido?
    \item ¿Qué rol recomienda para el adulto responsable dentro de la
          plataforma?
\end{enumerate}

\emph{Eje 2 -- reputación y diseño seguro.}

\begin{enumerate}
    \item ¿Qué errores evitaría en el diseño de un sistema de reputación
          bidireccional que involucra a menores?
    \item ¿Qué balance propone entre la protección del menor y el respeto a
          su privacidad al revisar sesiones o interacciones?
    \item ¿Qué recomendaciones daría para el manejo de denuncias y
          situaciones de conflicto entre tutores y familias?
\end{enumerate}

\subsection{Transcripción de la entrevista}

\completar{Transcripción organizada por ejes temáticos.}
        %!TEX root = ../main.tex
% APÉNDICE C — Detalles técnicos e implementación

\chapter{Detalles técnicos e implementación}\label{apx:tecnicos}

\begin{guia}
Todo lo que alguien necesita para reproducir los resultados:
hiperparámetros exactos, identificadores y versiones de los datos,
versiones de librerías, variables de entorno (sin secretos) y esquemas
de base de datos. Reproducibilidad total es un marcador de excelencia.
\end{guia}

\section{Configuración de los modelos}

\begin{table}[H]
    \centering
    \caption{Hiperparámetros del modelo final (ejemplo de estructura)}
    \label{tab:hiperparametros}
    \small
    \begin{tabular}{l l l}
        \toprule
        \tb{Parámetro} & \tb{Valor} & \tb{Criterio de selección} \\
        \midrule
        \completar{n\_estimators} & \completar{500} & \completar{búsqueda en grilla, validación cruzada 5 pliegues} \\
        \completar{max\_depth} & \completar{8} & \completar{...} \\
        \bottomrule
    \end{tabular}
\end{table}

\section{Datos utilizados}

\completar{Identificador, fuente, versión, fecha de descarga, licencia y
tamaño de cada conjunto de datos.}

\section{Arquitectura de la aplicación}

\subsection{Stack tecnológico}

\begin{lstlisting}[caption={Dependencias principales (ejemplo)}, label={lst:dependencias}]
python == 3.12
scikit-learn == 1.5.0
fastapi == 0.111.0
postgresql == 16
\end{lstlisting}

\subsection{Esquema de base de datos}

\begin{lstlisting}[language=SQL, caption={Definición de tablas (ejemplo de código SQL)}, label={lst:esquema}]
CREATE TABLE usuario (
    id          SERIAL PRIMARY KEY,
    correo      TEXT NOT NULL UNIQUE,
    creado_en   TIMESTAMPTZ NOT NULL DEFAULT now()
);
\end{lstlisting}

\subsection{Variables de entorno}

\completar{Nombre, propósito y ejemplo de valor de cada variable. Nunca
incluir secretos reales.}

        %!TEX root = ../main.tex
% APÉNDICE D — Registro de prompts utilizados con herramientas de IA
% Se referencia desde preliminares/declaracion-ia.tex (\ref{apx:prompts}).

\chapter{Prompts utilizados con herramientas de IA generativa}\label{apx:prompts}

\begin{guia}
Registro de los prompts más relevantes del trabajo, organizado por
capítulo y tema. Para cada uno: el comando \instruccionIA (qué se le pidió
a la herramienta) y su resultado \respuestaIA (lo que devolvió, ya
revisado y editado por el autor). Se incluyen solo los prompts cuyo
impacto en el documento fue significativo; el registro completo se
conserva en la carpeta compartida con la cátedra (política de uso de IA,
preliminares/declaracion-ia).
\end{guia}

\section{Criterios de registro}

\begin{itemize}
    \item \textbf{Trazabilidad}: cada entrada indica la herramienta y
          versión utilizada y la fecha aproximada.
    \item \textbf{Verificación}: toda cifra, cita o referencia sugerida
          por la herramienta se confirmó en la fuente original antes de
          incorporarla (regla 3 de la declaración de IA).
    \item \textbf{Autoría}: las salidas se transcriben solo como
          evidencia; el texto final del documento puede diferir por la
          revisión y adaptación del autor.
\end{itemize}

\section{Asistencia en el marco teórico y estado del arte}

\instruccionIA[1]{Resumí los argumentos principales de este paper sobre
evaluación de tutores y ranking de compatibilidad en plataformas de
tutoría en línea, en tres párrafos en español.}
\respuestaIA[1]{(resumen devuelto por la herramienta, verificado contra
el texto original)}

\instruccionIA[2]{Compará los modelos de negocio, comisiones y esquemas de
pago de Superprof, Preply, Wyzant y Tus Clases de manera tabular.}
\respuestaIA[2]{(tabla devuelta, contrastada con las tarifas públicas de
cada plataforma antes de usarse en \ref{tab:comparativa})}

\section{Asistencia en la justificación económica}

\instruccionIA[3]{Calculá el VAN a tasa 15\,\% y la TIR del flujo de caja
[Año 0 ... Año 5] de esta tabla de volúmenes y costos.}
\respuestaIA[3]{(cálculos devueltos; se replicaron de forma independiente
en hoja de cálculo antes de incluirlos en \ref{tab:flujo-caja} y en la
sección de indicadores)}

\section{Asistencia en redacción y código}

\instruccionIA[4]{Reescribí este párrafo en registro académico impersonal
sin cambiar su contenido.}
\respuestaIA[4]{(texto devuelto por la herramienta, luego editado)}

\instruccionIA[5]{Generá el scaffolding de una clase Java de Spring Boot
para el endpoint de reserva de una sesión, con validación de permisos por
rol.}
\respuestaIA[5]{(código devuelto, revisado y adaptado por el autor; las
pruebas se escribieron a mano)}
    \end{appendices}

    % anexos: material de terceros o externo (I, II...)
    \begin{anexos}
        %!TEX root = ../main.tex
% ANEXO I — Repositorios de código y material externo

\chapter{Repositorios de código fuente}\label{anx:repositorios}

\begin{guia}
Los anexos contienen material que no fue producido por el autor o que
es externo al documento: repositorios, datasets de terceros, normativa,
manuales del proveedor. Para cada repositorio: enlace, tecnología,
contenido y estado (activo, archivado).
\end{guia}

El código fuente de \sistema{} se aloja en un único repositorio público
de GitHub, organizado como monorepo.

\section{usal-tinku}

\prettygitlink{smercadocarbone/usal-tinku}

\textbf{Contenido.} El repositorio agrupa los tres componentes
desplegables del sistema, cada uno en su propia carpeta y con su propio
\texttt{README} de ejecución local:

\begin{itemize}
    \item \texttt{backend/}: \gls{api} REST en Java con Spring Boot,
          organizada como monolito modular.
    \item \texttt{matching-service/}: servicio de \gls{matching}
          semántico en Python, único proceso separado del monolito.
    \item \texttt{frontend/}: cliente \gls{pwa} en Next.js y React.
\end{itemize}

Se incluyen además la documentación de diseño (\texttt{docs/}), los
flujos de integración continua (\texttt{.github/workflows/}) y un
\texttt{docker-compose.yml} con la infraestructura local de desarrollo.

\textbf{Estado.} Activo, en desarrollo durante el cursado de la materia.

\completar{Licencia del repositorio (actualmente no declarada).}

        %!TEX root = ../main.tex
% ANEXO II — Modelo financiero
% Archivo: anexos/Presupuesto_y_Proyecciones_-_Proyecto_Tinku.xlsx

\chapter{Modelo financiero}\label{anx:modelo-financiero}

Los cálculos de la evaluación económica del
capítulo~\ref{cap:economica} se realizan en una planilla de cálculo con
parámetros editables, que se entrega junto con este documento en el
siguiente archivo:

\begin{center}
    \path{Presupuesto_y_Proyecciones_-_Proyecto_Tinku.xlsx}
\end{center}

Todos los montos están expresados en dólares estadounidenses.

\section{Estructura de la planilla}

\begin{table}[H]
    \centering
    \caption{Hojas del modelo financiero}
    \label{tab:hojas-modelo}
    \small
    \begin{tabularx}{\textwidth}{l L}
        \toprule
        \tb{Hoja} & \tb{Contenido} \\
        \midrule
        Inversión Inicial        & Inversión incremental del lanzamiento: horas y tarifa por perfil, dominio, asesoría legal y contingencia, más el costo hundido del \gls{mvp} a título informativo (tabla~\ref{tab:inversion}) \\
        Gastos Post-lanzamiento  & Costos fijos mensuales por año: infraestructura, servidor de aplicación, video, asistente \gls{llm}, honorarios contables, dominio y marketing (tabla~\ref{tab:costos-fijos}) \\
        Valor de Venta           & Comisión, precio por hora, duración promedio de las sesiones, ticket, costo de la pasarela, costo, precio y adopción del resumen opcional, y margen neto por sesión (tabla~\ref{tab:unit-economics}) \\
        Estimación de Demanda    & Alumnos activos por año y escenario, y sesiones anuales resultantes (tabla~\ref{tab:demanda}) \\
        VAN, TIR y Repago        & Tasa de descuento construida (tabla~\ref{tab:tasa}), flujo de fondos de los tres escenarios, \gls{van}, \gls{tir}, repago simple y descontado, y resumen comparativo (tablas~\ref{tab:flujo-caja} y~\ref{tab:indicadores-financieros}) \\
        Gráficos                 & Flujo neto y flujo acumulado por escenario, y \gls{van} comparado \\
        \bottomrule
    \end{tabularx}
\end{table}

\section{Supuestos editables}

Los valores de entrada están separados de los cálculos: al modificar
cualquiera de ellos se recalculan los tres escenarios. Los
principales son la comisión de plataforma, el precio por hora, la
proporción de sesiones de 30 minutos, el precio y la adopción del resumen opcional, las tarifas horarias por perfil,
los alumnos activos por año, las clases por alumno por mes y la tasa de
descuento. Cada supuesto documenta su origen en la columna de
observaciones o en un comentario de la celda.

Los indicadores se calculan con las funciones nativas \texttt{NPV} e
\texttt{IRR}, y los períodos de repago con la interpolación de la
ecuación~\eqref{eq:repago}.

    \end{anexos}

    % glosario, siglas y símbolos
    \glosarios

\end{document}
